\documentclass[twocolumn]{article}
\usepackage{times}
\usepackage{graphicx}
\usepackage{amsmath,amssymb,amsthm}
\usepackage{hyperref}
\usepackage{url}
\usepackage{microtype}
\usepackage{enumitem}
\usepackage{natbib}

\newcommand{\dtva}{\textsc{D-TVA}}
\newcommand{\tva}{\textsc{TVA}}
\newcommand{\tvv}{\textsc{TVV}}
\newcommand{\lr}{\textsc{LR}}

\newtheorem{definition}{Definition}

\title{%
  Dynamic Topological Void Analysis (D-TVA):\\
  Lifted Retrieval: Toward LLM-Conditioned\\
  Geometric Information Retrieval
}

\author{
  Pan, Kris \\
  Intel Corporation \\
  \texttt{kris.pan@intel.com}
}

\date{}

\begin{document}

\maketitle

\begin{abstract}
Dense contrastive encoders such as BGE-M3 provide a continuous geometric
substrate for information retrieval, but their InfoNCE training objective
enforces symmetry---making it \emph{mathematically impossible} to encode
asymmetric semantic relations such as hierarchical depth, entailment, or
causal direction~\cite{wang2020understanding}.
We address this limitation in the context of \emph{semantic drift detection}
over an evolving corpus.

We propose \textbf{Lifted Retrieval} (\lr{}), a paradigm in which:
(i)~a dense encoder provides a continuous geometric substrate,
(ii)~a chosen anchor restricts attention to a local neighborhood, and
(iii)~a large language model acts as a directional semantic operator
over that neighborhood, supplying the asymmetric structure the encoder
cannot represent.
Anchor conditioning is essential: it constrains LLM inference to
semantically hot regions, making the architecture deployable at corpus scale.

We instantiate \lr{} as \dtva{} (Dynamic Topological Void Analysis),
the third in a series of systems---following \tva{}~\cite{pan2026tva}
and \tvv{}~\cite{pan2026tvv}---that, in retrospect, all share this
architectural pattern.
In \dtva{}, the LLM operator supplies a hierarchy axis used to construct
anchor-conditioned entailment cones, measure drift polarity, and predict
research void evolution.

We plan to evaluate on the Linux Kernel Mailing List archive (1995--2026,
$>$1M messages), where RFC events provide explicit human-expert void
declarations with known outcomes.
We further introduce \textbf{LKHJB} (Linux Kernel Hierarchy Judgement
Benchmark), 200 expert-curated concept pairs for evaluating
LLM-as-hierarchy-oracle in specialized technical domains.
Experiments are ongoing; this preprint establishes the paradigm
and benchmark design.
\end{abstract}

%% ---------------------------------------------------------------
\section{Introduction}
\label{sec:intro}

Modern information retrieval systems built on dense sentence encoders
excel at topical proximity but are blind to \emph{direction}.
A query can retrieve documents near a concept, but the encoder cannot
distinguish whether a document is a generalization, a specialization,
or a sibling of that concept.
This blindness is not an artifact of insufficient training: it is a
consequence of the symmetric contrastive objective (InfoNCE) used to
train state-of-the-art encoders including BGE-M3, E5, and GTE.
By the alignment--uniformity framework of~\citet{wang2020understanding},
InfoNCE maximizes cosine similarity for positive pairs without any
asymmetric supervision signal.
No amount of scale or fine-tuning within this paradigm can produce a
direction-aware encoder.%
\footnote{The name ``Lifted Retrieval'' was assigned on 2026-05-22
during the design of \dtva{}, upon recognizing the recurrent
architectural pattern across three prior systems~\cite{pan2026tva,
pan2026tvv}. This preprint serves as the priority record for the term.}

We propose to resolve this asymmetry gap through \textbf{Lifted Retrieval}:
rather than modifying the encoder, we \emph{lift} the retrieval system
into a richer semantic space by introducing a large language model as a
directional operator, anchored to a local neighborhood to control cost.

\paragraph{Why this paradigm, why now.}
Three converging trends make this the right moment.
(i)~Contrastive encoders have reached diminishing returns on standard
benchmarks, indicating fundamental rather than incremental limitations.
(ii)~Open-weight LLMs (Llama-3.3-70B, Qwen-2.5-72B) now offer
inference at $<$\$0.001 per call, making per-neighborhood LLM operators
economically deployable.
(iii)~The symmetry impossibility~\cite{wang2020understanding} clarifies
why no scale-up of contrastive training will close the asymmetry gap.
\lr{} is not a method but a recognition of these convergent constraints.

\paragraph{Relation to prior work.}
In retrospect, our earlier work on void taxonomy~\cite{pan2026tva} and
void validation~\cite{pan2026tvv} both employed this pattern: dense
encoder for geometry, anchor for scope, LLM for semantic judgement.
Neither paper explicitly named the paradigm.
\dtva{} is the conscious naming moment.
A forthcoming framework paper will formalize the full operator algebra.

\paragraph{Contributions.}
\begin{enumerate}[leftmargin=*,itemsep=1pt]
\item \textbf{Lifted Retrieval paradigm} (Definition~\ref{def:lr}):
  formal definition of the encoder--anchor--LLM operator architecture.
\item \textbf{D-TVA}: an instantiation of \lr{} for anchor-conditioned
  semantic drift detection with Bures-Wasserstein drift metric and
  LLM-calibrated hierarchy axis.
\item \textbf{LKHJB benchmark}: 200 expert-curated Linux kernel
  hierarchy pairs with multi-model evaluation protocol, released
  as open data.
\item \textbf{Linux RFC evaluation (ongoing)}: event-driven prediction
  of RFC outcomes over 25 years of kernel development history.
\end{enumerate}

%% ---------------------------------------------------------------
\section{Related Work}
\label{sec:related}

\paragraph{Dense retrieval and its limits.}
State-of-the-art dense encoders (BGE-M3~\cite{bge-m3}, E5, GTE) achieve
high performance on symmetric relevance tasks but lack asymmetric
semantic structure by construction~\cite{wang2020understanding}.
Reranking with cross-encoders partially addresses this but at full-corpus
inference cost.

\paragraph{Hyperbolic embeddings for hierarchy.}
Poincar\'e embeddings~\cite{nickel2017poincare} and hyperbolic probing
methods~\cite{chen2021probing} show that transformer representations
contain latent hyperbolic structure.
However, extracting this structure from InfoNCE-trained encoders via
unsupervised lifting is unreliable: the Rayleigh quotient of contrastive
PCA is systematically biased toward sibling-class axes rather than the
abstraction depth axis.

\paragraph{LLM-augmented retrieval.}
Prior work uses LLMs for query expansion, reranking, or summarization
(RAG pipelines), but treats LLMs as supplementary tools rather than as
\emph{primary semantic operators} that resolve a structural limitation
of the encoder.
\lr{} provides a principled architectural foundation for this distinction.

%% ---------------------------------------------------------------
\section{Lifted Retrieval: Definition and Architecture}
\label{sec:lr}

\begin{definition}[Lifted Retrieval]
\label{def:lr}
A retrieval system instantiates the \emph{Lifted Retrieval} (\lr{})
paradigm if it composes three components:
\begin{enumerate}[leftmargin=*,itemsep=1pt]
\item A \textbf{dense geometric substrate}
  $\mathcal{E}: \text{text} \to \mathbb{R}^d$
  providing continuous topical proximity (typically a contrastive
  sentence encoder).
\item An \textbf{anchor} $C \in \mathbb{R}^d$ defining a local region
  of interest $\mathcal{N}_C \subset \mathbb{R}^d$, typically the
  $k$-nearest neighbors of $C$ under $\mathcal{E}$.
\item A \textbf{directional semantic operator}
  $\mathcal{T}: \mathcal{N}_C \to \mathcal{S}$, realized by a large
  language model, supplying asymmetric semantic structure $\mathcal{S}$
  that $\mathcal{E}$ cannot represent.
\end{enumerate}
The system is \emph{lifted} in the sense that $\mathcal{T}$ provides
capabilities beyond what $\mathcal{E}$'s symmetric training objective
admits.
\end{definition}

\paragraph{Scope and non-claims.}
We do not claim that LLMs are the only possible instantiation of
$\mathcal{T}$; any system providing asymmetric semantic structure
(e.g., knowledge graphs, symbolic reasoners) could in principle
serve as the operator.
We do not claim \lr{} is novel in its individual components---encoders,
anchors, and LLMs are widely used.
The contribution is the \emph{architectural composition} and the
explicit recognition that this composition resolves a \emph{fundamental}
representational limitation rather than an empirical one.
This preprint establishes the paradigm and benchmark protocol;
empirical validation is ongoing.

\begin{figure}[t]
\centering
\footnotesize
\begin{verbatim}
Corpus (docs over time)
  | encoder E (BGE-M3)
  v
Geometric Substrate    <-- cheap, indexable
  | + Anchor C
  v
Neighborhood N_C       <-- bounds LLM cost
  | + LLM operator T
  v
Directional Signal S   <-- hierarchy, entailment
  | + aggregation
  v
Output:
  D-TVA: drift/void
  TVA:   void taxonomy
  TVV:   void validation
\end{verbatim}
\caption{Lifted Retrieval architecture.
\dtva{}, \tva{}~\cite{pan2026tva}, \tvv{}~\cite{pan2026tvv}
are three operator instantiations.}
\label{fig:lr-arch}
\end{figure}

%% ---------------------------------------------------------------
\section{D-TVA: Lifted Retrieval for Semantic Drift}
\label{sec:dtva}

\dtva{} instantiates \lr{} with operator type
$\mathcal{T} = \text{hierarchy-score}$:
the LLM assigns an abstraction-depth score to each neighbor
$n \in \mathcal{N}_C$, enabling construction of an
\emph{anchor-conditioned hierarchy axis} $\hat{u}_h$ via ridge
regression in the anchor tangent space.

The full D-TVA pipeline is:
(1)~sphere log-map at anchor $C$;
(2)~generalized contrastive PCA (gcPCA) for anchor-distinctive
subspace $D^*$;
(3)~LLM-supervised hierarchy axis $\hat{u}_h$;
(4)~Bures-Wasserstein drift metric $dD_C(t_1,t_2)$ on the
Gaussian state $S_C(t) = \mathcal{N}(\mu_C(t), \Sigma_C(t))$;
(5)~drift polarity $\sigma_C = \operatorname{sign}(\langle
\Delta\mu_C, \hat{u}_h \rangle)$.

Mathematical details, implementation warnings, and failure modes
are documented in the companion technical report
\texttt{notes/anchor\_drift\_metric.md}.

%% ---------------------------------------------------------------
\section{LKHJB: Linux Kernel Hierarchy Judgement Benchmark}
\label{sec:lkhjb}

The \textbf{Linux Kernel Hierarchy Judgement Benchmark} (LKHJB) is a
set of 200 expert-curated (general, specific) concept pairs drawn from
8 kernel subsystems: scheduler, memory management, networking,
filesystem, block layer, locking, security, and tracing.

Each pair satisfies at least one of: specific \emph{is-a} general,
specific \emph{implements} general, specific \emph{is-a-mechanism-for}
general, or specific \emph{is-a-special-case-of} general.
Ambiguous sibling pairs (e.g., mutex vs.\ spinlock) are excluded.

The evaluation protocol uses pairwise LLM judgement with position
randomization (each pair evaluated in both orderings; correct only if
both orderings agree), reported as accuracy against the expert labels.

LKHJB will be released to Zenodo alongside the full paper.
Preliminary comparisons across unsupervised geometric baselines (gcPCA,
density gradient) and LLM oracles are in progress.

%% ---------------------------------------------------------------
\section{Open Questions}
\label{sec:open}

The \lr{} paradigm raises several questions for future work:
\begin{enumerate}[leftmargin=*,itemsep=1pt]
\item \textbf{Operator algebra.} What is the complete class of
  semantic structures expressible by autoregressive LLMs but not by
  symmetric contrastive encoders?
\item \textbf{Domain generalization.} Does the LLM-as-hierarchy-oracle
  transfer to domains where LLM pretraining coverage is lower (clinical,
  legal, proprietary)?
\item \textbf{Anchor selection.} What makes a good anchor, and can
  anchor quality be measured without human labeling?
\item \textbf{Operator cost vs.\ accuracy tradeoff.} What is the
  minimum LLM size that provides reliable directional semantics,
  and how does this interact with domain specificity?
\item \textbf{Dynamic anchor evolution.} How should the anchor itself
  be updated as the corpus evolves?
\end{enumerate}

%% ---------------------------------------------------------------
\section{Conclusion}
\label{sec:conclusion}

We introduced \textbf{Lifted Retrieval} (\lr{}), a paradigm that
resolves the fundamental asymmetry limitation of contrastive encoders
by composing them with an anchor-conditioned LLM operator.
The paradigm is instantiated as \dtva{} for semantic drift detection,
and is retrospectively recognized as the common pattern across three
prior systems.
The LKHJB benchmark and the Linux RFC evaluation are in progress.
We release this preprint as the priority record for the \lr{} paradigm.

\bibliographystyle{plainnat}
\bibliography{paper3}

\end{document}
